\appendix

% ==========================================================================
\section{Criterion scale: the cost of regressing on a percentile}
\label{app:uniform}

Conditional growth is published as a percentile rank. Suppose the underlying
growth construct $G$ is normally distributed and the screener relates to it
linearly. The published criterion is $P = \Phi(G)$, which is uniform on $(0,1)$.
Even a perfect monotone relationship between $G$ and $P$ is not a perfect linear
one, so regressing on $P$ loses information.

The size of the loss is the correlation between a standard normal variate and
its own probability integral transform. Integrating by parts with $u=\Phi(z)$
and $dv = z\varphi(z)\,dz$, so $v=-\varphi(z)$:

\[
\mathbb{E}[Z\,\Phi(Z)] = \Big[-\Phi(z)\varphi(z)\Big]_{-\infty}^{\infty}
  + \int_{-\infty}^{\infty}\varphi(z)^2\,dz
  = 0 + \frac{1}{2\pi}\int e^{-z^2}dz = \frac{1}{2\sqrt{\pi}}.
\]

Since $\mathbb{E}[Z]=0$ this is also $\operatorname{Cov}(Z,\Phi(Z))$, and
$\operatorname{Var}(\Phi(Z))=1/12$ because $\Phi(Z)$ is uniform. Hence

\[
\operatorname{corr}(Z,\Phi(Z))
 = \frac{1/(2\sqrt{\pi})}{1\cdot\sqrt{1/12}}
 = \frac{\sqrt{12}}{2\sqrt{\pi}} = \sqrt{\tfrac{3}{\pi}} = \UniformAttenR.
\]

The attainable $\Rsq$ is therefore multiplied by $3/\pi = \UniformAttenRsq$:
about \UniformAttenLossPct{} of explainable variance is discarded by using the
percentile directly. Applying $\Phi^{-1}$ to the conditional growth percentile
before analysis recovers it. The constant is verified numerically against
quadrature in \texttt{selftest.py}.

% ==========================================================================
\section{Stage 1: precision and power for a correlation}
\label{app:stage1}

For a bivariate normal sample, $z=\operatorname{artanh}(r)$ is approximately
normal with standard error $1/\sqrt{n-3}$. A two-sided test of
$H_0\!:\rho=0$ at level $\alpha$ with power $1-\beta$ against a true $\rho$
requires

\[
n \;=\; \left(\frac{z_{1-\alpha/2}+z_{1-\beta}}{\operatorname{artanh}\rho}\right)^{2} + 3 ,
\]

and the confidence interval is obtained by transforming
$\operatorname{artanh}(r) \pm z_{1-\alpha/2}/\sqrt{n-3}$ back through $\tanh$.
The routine reproduces Cohen's tabled value of $n=85$ for $\rho=0.30$ exactly;
at $\rho=0.50$ it returns 30 against Cohen's 28, the normal approximation being
mildly conservative at small $n$, which errs toward larger samples and is
therefore safe for planning.

\textbf{Why the report emphasises interval width over power.} A significance
test answers whether validity is distinguishable from zero. No admissions policy
turns on that. The decision-relevant quantity is the width of the interval,
because the cut score and the expected decision accuracy both depend on the
magnitude. This is why Section~\ref{sec:stage1} reports the sample needed for a
$\pm0.10$ interval (\NciTenAtThirtyFive{} children at $r=0.35$) alongside the
much smaller sample needed for significance (\NforThirtyFive{}).

\textbf{Multiplicity.} Two primary criteria are tested. The tables report the
Bonferroni-adjusted requirement ($\alpha=0.025$) as well as the unadjusted one.
Bonferroni is conservative here because the two criteria are positively
correlated; it is used because the correlation between them is unknown, and a
conservative adjustment is easier to defend than an estimated one.

\textbf{Clustering.} With clusters of size $m$ and intraclass correlation
$\mathrm{ICC}$, the variance of an estimate is inflated by the design effect
$1+(m-1)\mathrm{ICC}$, and the effective sample is $n$ divided by that factor.

% ==========================================================================
\section{Criterion contamination: bounds and the detection test}
\label{app:contam}

\subsection{How large a confounder would have to be}

Let $S$ be the sealed screener score, $P$ mastery pace, and $R$ a routing
variable. If $R$ is the only path from $S$ to $P$ other than the child's
ability, the contribution of that path to the observed correlation is
$\rho_{SR}\rho_{RP}$. To account for an observed $r_{SP}$ in full, the product
of the two legs must equal $r_{SP}$; the symmetric case sets both legs to
$\sqrt{r_{SP}}$.

\begin{center}
\small
\begin{tabular}{@{}cccc@{}}
\toprule
\textbf{Observed $r_{SP}$} & \textbf{Both legs equal} &
\textbf{$\rho_{RP}$ if $\rho_{SR}{=}0.30$} &
\textbf{$\rho_{RP}$ if $\rho_{SR}{=}0.50$} \\
\midrule
\tabinput{generated/tab_contamination_sensitivity.tex}
\bottomrule
\end{tabular}
\end{center}

At an observed validity of $0.35$, a confounder would need both legs at
\ConfoundSymThirtyFive{}, or --- if routing correlates only $0.30$ with the
screener --- a routing-to-pace association above $1$, which is impossible. The
reading is that \emph{weak} routing dependence on the screener cannot explain a
moderate validity; only strong dependence can. That is a specific, checkable
claim rather than an unfalsifiable worry.

This is the logic of sensitivity analysis for unmeasured confounding: rather
than asserting no confounding, state how strong confounding would have to be,
and let the reader judge whether that is plausible.

\subsection{The detection test}

Contamination inflates $\operatorname{corr}(S,P)$ but not
$\operatorname{corr}(S,M)$ where $M$ is MAP conditional growth, because MAP is
outside the platform's routing. The diagnostic is therefore the difference
between two correlations that share the variable $S$, tested by the
Meng--Rosenthal--Rubin statistic:

\[
Z=(z_1-z_2)\sqrt{\frac{n-3}{2(1-r_{PM})\,h}},
\qquad
h=\frac{1-f\,\bar r^{2}}{1-\bar r^{2}},
\quad
f=\min\!\left(\frac{1-r_{PM}}{2(1-\bar r^{2})},\,1\right),
\]

with $z_i=\operatorname{artanh}(r_i)$ and $\bar r^{2}=(r_1^{2}+r_2^{2})/2$.
Required samples:

\begin{center}
\small
\begin{tabular}{@{}ccccc@{}}
\toprule
\textbf{$r$(screener, pace)} & \textbf{$r$(screener, MAP)} &
\multicolumn{3}{c}{\textbf{$n$ for 80\% power, by $r$(pace, MAP)}} \\
\cmidrule(lr){3-5}
 & & $0.30$ & $0.50$ & $0.70$ \\
\midrule
\tabinput{generated/tab_contamination_power.tex}
\bottomrule
\end{tabular}
\end{center}

Note the direction of the dependence on $r_{PM}$: the more strongly pace and MAP
agree, the \emph{easier} the discrepancy is to detect, because the comparison
becomes a within-child contrast. Detecting a gap between $0.45$ and $0.30$ at
$r_{PM}=0.50$ requires \NcontamCheckMid{} children; the pooled cohort of
\Nbothcohorts{} reaches power \PowerBothContam{}. GT Anywhere alone
(\NanywhereCohort{} children) reaches \PowerAnywhereContam{}, so the check does
not depend on pooling the two cohorts --- which matters, because pooling a
virtual and a physical cohort introduces exactly the clustering discussed in
Appendix~\ref{app:stage1}.

\subsection{Interpreting the outcome}

\begin{itemize}[leftmargin=1.4em,itemsep=0.15em]
\item \textbf{No detectable gap.} Pace and MAP agree about the screener.
Contamination is bounded below the detectable difference, and the primary
criterion can be used with its full validity estimate.
\item \textbf{Pace materially exceeds MAP.} Consistent with routing
contamination, and also with pace being a more proximal and less noisy measure
of what the screener predicts. These are not distinguishable by this test alone;
the exposure-adjusted analysis in Section~\ref{sec:contamination} is what
separates them. Until it does, the MAP figure should be treated as the headline.
\item \textbf{MAP exceeds pace.} Not contamination. Most likely the pace measure
is noisy, or its denominator is mis-specified.
\end{itemize}

\boundary{This test bounds contamination; it does not eliminate it. A routing
mechanism that happened to affect pace and MAP equally would be invisible to it.
That is implausible given MAP's externality but is not excluded, and it is the
reason the routing-input audit in Section~\ref{sec:contamination} is listed
ahead of the statistical checks.}

% ==========================================================================
\section{Equating, linking and concordance}
\label{app:equating}

The measurement literature distinguishes several kinds of relationship between
scores on two tests, in decreasing order of strength.

\textbf{Equating} produces interchangeable scores. It requires that the two
tests measure the same construct, at comparable reliability, with the same
intended use, and that the conversion be invariant across subpopulations. Only
then may a score on one be substituted for a score on the other without
qualification.

\textbf{Calibration and projection} are weaker forms, applicable when tests
measure the same construct at different reliabilities or difficulties.

\textbf{Concordance} is a purely statistical correspondence between scores on
tests measuring \emph{different} constructs. It is population-dependent: the
concordance table derived in one group need not hold in another. It supports
statements of the form ``in this population, children at this screener score
tended to obtain that CogAT score,'' and nothing stronger.

\textbf{Why this matters for GT.} A concordance table between the screener and
CogAT would be straightforward to produce once both are administered, and it
would be actively harmful if labelled an equating. It would invite the inference
that the screener can substitute for CogAT in contexts where CogAT is required
--- district gifted identification, for instance --- which the evidence would
not support and which would expose GT to a defensible complaint.

More importantly, concordance answers a question GT does not need answered.
Knowing that a screener score of $X$ corresponds to a CogAT SAS of $Y$ tells GT
nothing about whether the screener should replace CogAT. That decision rests on
which instrument better predicts the criterion, and on whether the screener adds
information CogAT lacks. Both are incremental-validity questions, which is why
Section~\ref{sec:stage2} is framed as it is.

\textbf{The invariance question is separate and still required.} Before
operational use, the screener must be checked for differential item functioning
and for measurement invariance across the groups it will be used on, and for
differential prediction --- whether a given score forecasts the same criterion
value in every group. A test can be free of item bias and still predict
differently across groups; these are distinct properties requiring distinct
evidence.

% ==========================================================================
\section{The increment manufactured by unreliability}
\label{app:spurious}

\subsection{Setup}

Let $T$ be a latent ability, standardised. Let the criterion $Y$ be standardised
with $\operatorname{corr}(T,Y)=\rho$. Let two instruments be congeneric
indicators of $T$ with reliabilities $r_1$ and $r_2$:

\[
X_1=\sqrt{r_1}\,T+\sqrt{1-r_1}\,e_1,
\qquad
X_2=\sqrt{r_2}\,T+\sqrt{1-r_2}\,e_2,
\]

with $e_1,e_2$ independent of each other, of $T$, and of $Y$ given $T$. By
construction $X_2$ carries \emph{no information about $Y$ beyond $T$}: it is a
second thermometer for the same temperature. Write

\[
a=\operatorname{corr}(X_1,Y)=\sqrt{r_1}\,\rho,\quad
b=\operatorname{corr}(X_2,Y)=\sqrt{r_2}\,\rho,\quad
c=\operatorname{corr}(X_1,X_2)=\sqrt{r_1r_2}.
\]

\subsection{Derivation}

$R^2$ from $X_1$ alone is $a^2=r_1\rho^2$. For two predictors,

\[
R^2_{\text{full}}=\frac{a^2+b^2-2abc}{1-c^2}.
\]

Substituting, $a^2+b^2 = \rho^2(r_1+r_2)$ and
$2abc = 2\rho^2\sqrt{r_1r_2}\sqrt{r_1r_2}=2\rho^2 r_1r_2$, so the numerator is
$\rho^2(r_1+r_2-2r_1r_2)$ and

\begin{align}
\dR &= \frac{\rho^2(r_1+r_2-2r_1r_2)}{1-r_1r_2}-r_1\rho^2
     = \rho^2\,\frac{r_1+r_2-2r_1r_2-r_1(1-r_1r_2)}{1-r_1r_2} \notag\\[0.3em]
    &= \rho^2\,\frac{r_2-2r_1r_2+r_1^2r_2}{1-r_1r_2}
     = \boxed{\;\rho^{2}\,\frac{r_2\,(1-r_1)^{2}}{1-r_1r_2}\;}
\end{align}

The closed form is checked in \texttt{selftest.py} against $R^2$ computed
directly from the $3\times3$ correlation matrix, at several parameter
combinations, to machine precision.

\subsection{Reading it}

The expression behaves as it should. It vanishes when $r_1=1$: a perfectly
measured incumbent leaves nothing for a redundant challenger to pick up. It
grows with $\rho^2$, so the artefact is largest exactly when the construct
matters most. It grows with $r_2$, so a \emph{more} reliable redundant
challenger produces a \emph{larger} spurious increment --- which is the
counterintuitive part, and the reason a well-built screener is more likely to
produce a misleading positive than a poorly built one.

Substituting GT's figures --- comparator reliability $\RelCompLow$--$\RelCompHigh$,
a screener reliability target of $\RelScreener$, and a criterion correlation
between the Sackett-revised operational value $\RhoCogAToperational$ and the
observed meta-analytic value $\RhoCogATobserved$ --- gives the table in
Section~\ref{sec:floor}: an artefact between $\SpuriousLow$ and $\SpuriousHigh$,
with $\SpuriousMid$ as the central case.

\subsection{Consequence for the test}

At the central case, the conventional test against $\dR=0$ reaches 80\% power to
declare the \emph{artefact} significant at $n=\NdetectSpurious$; at the
pessimistic case, at $n=\NdetectSpuriousHigh$. Since Stage~2 will plausibly run
at a few hundred children, a naive analysis is at material risk of reporting a
significant increment that is entirely an artefact of CogAT's measurement error.
Testing against the floor instead raises the sample needed for a real
five-point increment from $\NincrFiveMid$ to $\NshiftedFive$.

\assumption{The floor depends on $r_2$, the screener's own reliability, which is
currently a design target of \RelScreener{} rather than a measurement. It also
assumes the two instruments are congeneric indicators of a single latent trait,
which is the worst case for GT --- it is the assumption under which the screener
adds nothing. If the screener genuinely measures a partly distinct construct,
the true floor is lower and the test is conservative.}

% ==========================================================================
\section{Decision accuracy and the base-rate arithmetic}
\label{app:base}

\subsection{The $2\times2$ table is implied, not chosen}

Sensitivity and specificity are often treated as properties of an instrument.
For a continuous score cut at a threshold they are not: they follow from the
validity, the selection ratio and the base rate. Model the screener $X$ and the
criterion $Y$ as standard bivariate normal with correlation $\rho$. Select
children with $X>z_x$ where $z_x=\Phi^{-1}(1-\mathrm{SR})$; define ``excels'' as
$Y>z_y$ where $z_y=\Phi^{-1}(1-\mathrm{BR})$. Then

\[
\mathrm{TP}=\Pr(X>z_x,\,Y>z_y),\quad
\mathrm{FP}=\mathrm{SR}-\mathrm{TP},\quad
\mathrm{FN}=\mathrm{BR}-\mathrm{TP},
\]
\[
\mathrm{TN}=1-\mathrm{SR}-\mathrm{FN},\quad
\text{Se}=\frac{\mathrm{TP}}{\mathrm{BR}},\quad
\text{Sp}=\frac{\mathrm{TN}}{1-\mathrm{BR}},\quad
\text{PPV}=\frac{\mathrm{TP}}{\mathrm{SR}} .
\]

The joint probability is a bivariate normal orthant, computed here by
one-dimensional quadrature and validated against Sheppard's theorem, against
Plackett's identity, and against the published Taylor--Russell tables.

A structural consequence worth noting: sensitivity is bounded above by
$\mathrm{SR}/\mathrm{BR}$. If GT selects the top $10\%$ but $20\%$ of children
meet the criterion, sensitivity cannot exceed $0.50$ however good the instrument
is, because there are not enough places. The low sensitivity figures in
Section~\ref{sec:decision} are largely this, not instrument failure.

\subsection{The odds form, as a cross-check}

The same arithmetic in likelihood-ratio form makes the base-rate dependence
transparent:

\[
\underbrace{\frac{\text{PPV}}{1-\text{PPV}}}_{\text{posterior odds}}
=\underbrace{\frac{\mathrm{BR}}{1-\mathrm{BR}}}_{\text{prior odds}}
\times
\underbrace{\frac{\text{Se}}{1-\text{Sp}}}_{\text{likelihood ratio}} .
\]

The instrument contributes only the likelihood ratio. At Se $=0.90$ and
Sp $=0.95$ the likelihood ratio is $18$, which is strong. At a base rate of
$2\%$ the prior odds are $1{:}49$, so the posterior odds are $18{:}49$ --- a
PPV of \PPVtwoPercentGood{}. \textbf{A likelihood ratio of 18 is an excellent
diagnostic and still leaves most flagged children below the bar}, because the
prior was so unfavourable. No improvement in the instrument fixes this; only
redefining the criterion, or accepting that the score surfaces candidates rather
than identifies successes, does.

\subsection{Decision consistency and accuracy}

Model observed score $X=T+E$ with $\operatorname{Var}(X)=1$ and
$\operatorname{Var}(T)=r_{xx}$. Two parallel administrations correlate $r_{xx}$,
so consistency at a cut $c$ is
$\Pr(X_1>c,X_2>c)+\Pr(X_1<c,X_2<c)$ under a bivariate normal with correlation
$r_{xx}$. Accuracy compares the observed classification to the true-score
classification at the same raw cut, using
$\operatorname{corr}(X,T)=\sqrt{r_{xx}}$ and the cut $c/\sqrt{r_{xx}}$ on the
standardised true score.

At a median cut this has the exact closed form
$\tfrac12+\arcsin(r_{xx})/\pi$, which \texttt{selftest.py} verifies to nine
decimal places. These are model-based figures assuming normality and classical
error; the operational estimate for a real test would use the
Livingston--Lewis procedure on the actual score distribution. The modelled
figures are adequate for design, and the report labels them as modelled.

% ==========================================================================
\section{Range restriction}
\label{app:range}

\subsection{Thorndike Case II}

Under direct selection on the predictor with unrestricted SD $S$ and restricted
SD $s$, write $U=S/s$. The observed correlation in the restricted group relates
to the unrestricted one by

\[
r_{\text{restricted}}
= \frac{\rho/U}{\sqrt{1-\rho^{2}+\rho^{2}/U^{2}}},
\qquad
\rho
= \frac{r_{\text{restricted}}\,U}{\sqrt{1-r_{\text{restricted}}^{2}+r_{\text{restricted}}^{2}U^{2}}} .
\]

For top-fraction selection with ratio $p$ and cut $c=\Phi^{-1}(1-p)$, the
truncated variance follows from the standard normal moments. With
$\lambda=\varphi(c)/p$ the inverse Mills ratio,

\[
\operatorname{Var}(X\mid X>c) = 1+\lambda c-\lambda^{2} .
\]

At $p=0.10$: $c=1.2816$, $\lambda=1.7550$, and the SD ratio is
$s/S=\SDratioTen$. The closed form is checked against numerical quadrature of
the truncated density in \texttt{selftest.py}, and the attenuation and
correction functions are checked to invert one another.

\subsection{What GT avoids}

A true validity of $0.40$ observed under top-$10\%$ selection appears as
$\RestrictedTenFromForty$, an attenuation of \AttenuationTenPct{}. The sample
needed to detect it at 80\% power rises from $\NunrestrictedForty$ to
$\NrestrictedForty$ --- a factor of \NmultiplierTen{}. At top-$20\%$ selection
the factor is still \NmultiplierTwenty{}.

The deeper benefit is that no correction is applied at all. Corrections require
the unrestricted SD as an input, assume linearity and homoscedasticity across a
range that was by definition never observed, and have been shown to
systematically overcorrect: the operational validity of cognitive ability was
revised downward from roughly $0.51$ to $0.31$ once the overcorrection was
identified. \textbf{An uncorrected estimate from an unrestricted sample is not
merely more efficient --- it is a different and stronger class of evidence,
because it does not depend on the correction machinery that produced that
error.}

\subsection{Indirect restriction, if it appears}

If selection operates on a composite that includes the screener rather than on
the screener alone, Case~II no longer applies and Case~III is required, which
additionally needs the correlations between the screener, the composite and the
criterion. Case~III corrections are substantially less stable. The mitigation is
procedural rather than statistical: keep the selection rule explicit and
recorded. An explicit rule is correctable; an unrecorded judgement is not.

% ==========================================================================
\section{Block length and the subscore-value bar}
\label{app:block}

\subsection{From recovery to reliability}

The repository reports \emph{recovery} $r$: the correlation between the
estimated learning rate and the injected true value in simulation. If the
estimate is the truth plus independent error, then
$\operatorname{corr}(\hat\lambda,\lambda)=\sqrt{\text{reliability}}$, so the
implied reliability of the readout is $r^{2}$. At the administered
\BlockAdministered{} trials, recovery of \RecoveryThirty{} implies a reliability
of \RelLRthirty{}; at 60 trials, \RecoverySixty{} implies \RelLRsixty{}.

\subsection{The subscore-value bar}

A subscore earns separate reporting only when it predicts its own true score
better than the total does --- the proportional-reduction-in-mean-squared-error
criterion. Writing $r_s$ for the subscore reliability, $r_x$ for the total's,
and $\rho$ for the correlation between their true scores, the condition is
$r_s > \rho^{2}r_x$, so the subscore clears the bar only when

\[
\rho \;<\; \sqrt{r_s/r_x}.
\]

At \BlockAdministered{} trials this ceiling is \HabermanThirty{}: the
learning-rate block earns separate reporting only if learning rate and standing
ability are correlated below \HabermanThirty{} in truth. At 60 trials the
ceiling relaxes to \HabermanSixty{}.

\subsection{The hard ceiling on incremental validity}

Separately, there is an upper bound on how much a subscore can add to prediction
of an \emph{external} criterion: the incremental validity is no greater than the
reliability of the residual from linear prediction of the subscore by the total.
With $S$ and $X$ standardised, $\operatorname{corr}(S,X)=c=\rho\sqrt{r_sr_x}$,
the residual is $e=S-cX$ with $\operatorname{Var}(e)=1-c^{2}$. Its true part is
$T_e=T_s-cT_x$, and since $\operatorname{Var}(T_s)=r_s$,
$\operatorname{Var}(T_x)=r_x$ and $\operatorname{Cov}(T_s,T_x)=c$,

\[
\operatorname{Var}(T_e)=r_s-2c^{2}+c^{2}r_x=r_s\big[1-\rho^{2}r_x(2-r_x)\big],
\]
\[
\dR_{\max}=\operatorname{rel}(e)
=\frac{r_s\big[1-\rho^{2}r_x(2-r_x)\big]}{1-\rho^{2}r_sr_x} .
\]

At $\rho=0$ this reduces to $r_s$: an orthogonal subscore can add at most its
own reliability, which at \BlockAdministered{} trials is \RelLRthirty{}. At
$\rho=\RhoLRstanding$ the ceiling falls to \HabCeilThirty{}.

\subsection{Sensitivity to the unmeasured correlation}

Both bars --- the subscore-value bar $\rho < \sqrt{r_s/r_x}$ and the incremental
ceiling above --- are governed by $\rho$, the true correlation between learning
rate and standing ability, which has never been measured. The sensitivity is
tabulated in Section~\ref{sec:twophase}. Two features of it deserve emphasis
here, because both are easy to miss from the main text.

First, the failure point at the administered block length is $\rho \approx
\HabermanThirty{}$, and the working assumption of \RhoLRstanding{} sits only
\HabermanMarginThirty{} below it. The component does not survive to
\RhoLRcollapse{} and then fail; it fails at $0.50$, which is an unremarkable
correlation between two cognitive constructs measured on the same children by
the same instrument on the same day.

Second, the two bars fail in different ways and the difference matters for how a
negative result should be reported. Exceeding the subscore-value bar means the
subscore should not be reported \emph{at all}, because the total score predicts
it better than it predicts itself. Falling below the incremental ceiling means
only that the subscore's contribution is bounded and small. The first is a
prohibition; the second is a magnitude. At \RhoLRcollapse{} the administered
block \HabermanCollapseVerdict{} the first, which is the stronger and more
damaging outcome.

\subsection{Attenuation, and the decisive comparison}

The ceiling is an upper bound. The \emph{expected} increment is governed by
attenuation: a predictor with reliability $r_s$ realises only $r_s$ of a true
$\dR$. A genuinely real increment of \TrueDeltaAssumed{} is realised as
$\RelLRthirty\times\TrueDeltaAssumed=\ObsDeltaThirty$ at \BlockAdministered{}
trials.

The companion simulation measured the realised increment directly, against the
external MAP anchor, and obtained \LRdeltaThirty{} at \BlockAdministered{}
trials and \LRdeltaSixty{} at 60. The agreement with the attenuation prediction
at \BlockAdministered{} trials is close, which is mild corroboration that the
attenuation model is describing the right mechanism.

\keyfinding{\textbf{Whether the measured increment clears the artefact floor
depends on an assumption, not on the measurement.} At the working comparator
reliability of \RelComparator{}, the floor is \SpuriousMid{} and the
\BlockAdministered{}-trial increment of \LRdeltaThirty{} sits
\LRthirtyVsFloor{} it. If the criterion correlation is at the upper end of its
plausible range the floor rises to \SpuriousHigh{} and the same increment sits
\LRthirtyVsFloorHigh{} it. If the comparator's reliability is \RelCompLow{}
rather than \RelComparator{}, the floor reaches \SpuriousWorst{} and the
increment is not close.
\smallskip

\noindent The honest summary is that the learning-rate component's contribution
at the administered block length is \emph{the same order of magnitude as the
artefact a second fallible measure of the same construct produces for free}. It
cannot be separated from that artefact without first pinning down the
comparator's reliability in GT's own sample --- which is a measurement GT can
make, and should, before running this comparison at all.}

At 60 trials the measured increment rises to \LRdeltaSixty{}, \LRsixtyVsFloor{}
the working floor, and detectable at \NlrSixty{} children rather than
\NlrThirty{}. Sixty trials is therefore the minimum plausible length for the
covariate claim --- necessary, and still not sufficient for an absolute claim,
because the contamination floor falls only \LRfloorDropPct{} across that
doubling while the random error falls \LRseDropPct{}
(Section~\ref{sec:covariate-vs-absolute}).

\assumption{The correlation between learning rate and standing ability,
$\rho=\RhoLRstanding$ in these tables, is an illustrative value. It has not been
measured and it drives both the subscore-value ceiling and the incremental
ceiling. Measuring it is the main reason to keep administering the block during
validation even though its output is not reportable.}

% ==========================================================================
\section{ROC precision, paired comparison, and the cost of the frame}
\label{app:auc}

This appendix supports Section~\ref{sec:feasibility}. It derives the precision
of an AUC estimate, the precision of the \emph{difference} between two AUCs
measured on the same children, and the translation between the AUC frame and
the variance frame used everywhere else in this report.

\subsection*{Precision of a single AUC}

For an AUC of $A$ estimated on $n_1$ positive and $n_0$ negative cases,
Hanley and McNeil (1982) give
\[
  \operatorname{Var}(A) = \frac{A(1-A) + (n_1-1)(Q_1 - A^2) + (n_0-1)(Q_2 - A^2)}
                               {n_1 n_0},
\]
\[
  Q_1 = \frac{A}{2-A},\qquad Q_2 = \frac{2A^2}{1+A},
\]
where $Q_1$ and $Q_2$ are the probabilities that two positives both outrank a
negative, and that a positive outranks two negatives. The expression is exact
under an exponential model for the score distributions and is used here as an
approximation.

Two checks are run against this in \path{scripts/selftest.py}. At $A=0.5$ it must
reduce exactly to the null variance of the Mann--Whitney statistic,
$\operatorname{Var} = (N+1)/(12\,n_1 n_0)$, which it does to machine precision;
and at $A=1$ it must vanish, which it does. At $n=\NcampusCohort{}$ with a
\PoolBaseRate{} base rate it returns a half-width of \AnalyticHalfwidth{},
against \MeasuredHalfwidth{} measured by subsampling the simulated cohort. The
agreement is close enough to treat either as the working figure.

\subsection*{Why the paired difference is far more precise}

The relevant quantity is not the width of either marginal interval but the
precision of $A_1 - A_2$ when both are computed on the \emph{same} children:
\[
  \operatorname{Var}(A_1 - A_2)
   = \operatorname{Var}(A_1) + \operatorname{Var}(A_2)
     - 2\,r_{\mathrm{AUC}}\sqrt{\operatorname{Var}(A_1)\operatorname{Var}(A_2)},
\]
with $r_{\mathrm{AUC}}$ the correlation between the two AUC statistics, itself
induced by the correlation between the two instruments' scores within the
positive and negative groups (Hanley \& McNeil, 1983). Because both instruments
measure overlapping reasoning constructs, $r_{\mathrm{AUC}}$ is large and the
shared sampling error cancels.

This is the reason Section~\ref{sec:feasibility} does not simply divide the
single-AUC half-width by the effect. Comparing two marginal intervals and
concluding ``no difference'' where they overlap is a well-documented error
(Gelman \& Stern, 2006). The correct paired test is several times more
efficient --- and still, at $r_{\mathrm{AUC}}=0.9$, requires \NsupHi{} children
for $80\%$ power against an effect of \AUCdiff{}.

$r_{\mathrm{AUC}}$ is not known for this pair of instruments and is treated as a
sensitivity parameter throughout, spanning $0.5$ to $0.9$. It can be estimated
directly once any sample has both scores, and doing so early is worthwhile: it
moves the required sample by a factor of five.

\subsection*{Non-inferiority}

A superiority test asks whether the lower bound on $A_1 - A_2$ exceeds zero. A
non-inferiority test asks whether it exceeds $-\delta$ for a pre-specified
margin $\delta$. With a true difference $\Delta$ the required sample follows from
\[
  \frac{\Delta + \delta}{\operatorname{SE}(A_1-A_2)} \;\ge\; z_{1-\alpha} + z_{1-\beta},
\]
one-sided at $\alpha = 0.025$. At $\delta = \NImargin{}$ this gives
\NnonInfHi{} children at $r_{\mathrm{AUC}}=0.9$ against \NsupHi{} for
superiority. The margin must be fixed before the data are seen and justified as
a difference GT would genuinely tolerate in exchange for whatever the screener
offers over CogAT --- cost, time, breadth, or accessibility. A margin chosen
afterwards to make the result come out is the standard abuse of this design and
is easy for a reviewer to detect from the analysis timestamp.

\subsection*{Translating between AUC and $R^2$}

The rest of this report works in correlations and variance increments; the
companion simulation reports AUCs. The bridge, for a continuous predictor $X$
and a latent criterion $Y$ that are bivariate normal with correlation $\rho$,
with the binary outcome defined as $Y$ exceeding the $1-p$ quantile $c$:
\[
  \mathrm{AUC} = \int_{-\infty}^{\infty} f_{+}(x)\,F_{-}(x)\,\mathrm{d}x,
  \qquad
  f_{+}(x) = \frac{\phi(x)\,\bar\Phi\!\big(\tfrac{c-\rho x}{\sqrt{1-\rho^2}}\big)}{p},
\]
and $f_{-}$ the same with $\Phi$ in place of $\bar\Phi$ and $1-p$ in the
denominator. This is evaluated by quadrature in \path{scripts/statlib.py} and
inverted by bisection.

It is validated against an independent reduction. Writing $U = X_i - X_j$, the
event $\{X_i > X_j,\ Y_i > c,\ Y_j \le c\}$ is an orthant probability of the
trivariate normal $(U/\sqrt2,\ Y_i,\ -Y_j)$ with correlations
$(\rho/\sqrt2,\ \rho/\sqrt2,\ 0)$; conditioning on the first coordinate reduces
it to a one-dimensional integral over the bivariate orthant routine already
validated against Sheppard's theorem and Plackett's identity. The two paths
agree to $10^{-6}$.

Applying it to the simulated results: an AUC of \AUCinstr{} at a \PoolBaseRate{}
base rate corresponds to $\rho = \RhoInstr$, and \AUCcomp{} to
$\rho = \RhoComp$. The first of those is worth noting --- the validity of
$0.35$ assumed throughout this report as a working value, chosen before the
simulation was run, lands within $0.01$ of what the simulation produced.

\subsection*{What dichotomising the criterion costs}

Computing an AUC requires a binary criterion. Mastery pace is continuous, so
using AUC means dichotomising it. For a bivariate normal pair split at base
rate $p$, the correlation with the dichotomised criterion is
\[
  r_{\mathrm{pb}} = \rho\,\frac{\phi(c)}{\sqrt{p(1-p)}},
\]
whose multiplier is at most $\sqrt{2/\pi} = 0.798$ at a median split and
\DichFactor{} at $p = \PoolBaseRate{}$. Squaring, about $42\%$ of the explained
variance is discarded --- information the study paid to collect and then threw
away at the analysis step (MacCallum et al., 2002). In sample terms the penalty
is a factor of \DichPenalty{}, and the further step to a paired AUC comparison
costs a factor of \AUCpenalty{} against the continuous variance increment.

The recommendation that follows is narrow and worth stating precisely: compute
AUC, report AUC, and put AUC in the slide deck, because it is the statistic a
non-technical audience can interpret. Do not \emph{test} on it, and do not
\emph{size the study} on it.

\subsection*{The asymmetry between the two increments}

Finally, the two directions of the incremental question are not the same
question and do not have the same answer. In the simulation:
\[
  R^2_{\text{comparator}} = \RtwoComp, \qquad
  R^2_{\text{screener}} = \RtwoInstr, \qquad
  R^2_{\text{both}} = \RtwoBoth.
\]
Adding the screener to the comparator gains \DeltaAddInstr{}; adding the
comparator to the screener gains \DeltaAddComp{}, which falls
\DeltaAddCompVsFloor{} the artefact floor of Appendix~\ref{app:spurious}. A
result of that shape does not license ``the screener is better.'' It licenses
``the comparator adds nothing measurable once the screener is in the model,''
which is a claim about redundancy rather than superiority, and is the claim
Section~\ref{sec:feasibility} recommends GT pursue.

\section{Reproducibility}
\label{app:repro}

\subsection{How the numbers were produced}

Every numeric value in this report is generated by
\texttt{scripts/compute.py} and written into \texttt{generated/} as \LaTeX{}
macros and table bodies, which the document then includes. No number is typed
into the prose by hand. To regenerate:

\begin{center}
\texttt{cd scripts \&\& python3 selftest.py \&\& python3 compute.py}
\end{center}

\texttt{scripts/statlib.py} implements the statistical routines using only the
Python standard library --- no NumPy, no SciPy --- so the results reproduce on a
bare Python 3.8+ interpreter. The routines include the regularised incomplete
beta function, the central and noncentral $F$ distributions, bivariate normal
orthant probabilities, and the power and sample-size functions built on them.

\subsection{Validation}

\texttt{scripts/selftest.py} checks those routines against external reference
values and internal consistency identities, and exits non-zero on any failure.
It is a prerequisite of the build. The external checks include:

\begin{itemize}[leftmargin=1.4em,itemsep=0.1em]
\item central $F$ quantiles against published tables at three
$(\mathrm{df}_1,\mathrm{df}_2,\alpha)$ combinations;
\item two G*Power worked examples for the $\Rsq$-increment test, reproduced
exactly ($n=77$ for $f^2=0.15$ with $q=3$; $n=395$ for $f^2=0.02$ with $q=1$);
\item Cohen's tabled sample size for a correlation ($n=85$ at $\rho=0.30$);
\item Taylor--Russell tabled success ratios at two parameter combinations;
\item Sheppard's theorem for the bivariate normal at the median cut, across five
correlations;
\item Plackett's identity for off-centre orthant probabilities;
\item the Hanley--McNeil AUC variance reduced to the Mann--Whitney null variance
$(N+1)/(12 n_1 n_0)$ at $A=0.5$, exactly, at three sample splits;
\item the AUC/correlation bridge against an independent trivariate-orthant
reduction of the same probability, agreeing to $10^{-6}$
(Appendix~\ref{app:auc}).
\end{itemize}

Internal checks confirm that the range-restriction correction inverts the
attenuation exactly, that the closed-form spurious increment matches $\Rsq$
computed from the correlation matrix, that the $\sqrt{3/\pi}$ constant matches
quadrature, that a zero-validity predictor selects exactly at the base rate, and
that a zero-reliability test agrees with itself only at chance.

\subsection{Building the document}

\texttt{make} regenerates the numbers and builds \texttt{main.pdf}. If no local
\LaTeX{} installation is present, \texttt{make docker} builds the same PDF in a
TeX~Live container. \texttt{make check} runs the self-test and a consistency
check that every macro referenced in the sources is defined in
\texttt{generated/numbers.tex}. That check catches three things the build
either misses or reports late: a stale \texttt{generated/} directory, a macro
that is computed but no longer referenced anywhere (usually a number left behind
by a rewritten passage), and a passage calling a macro that no longer exists.

\subsection{Figures}

The six figures are produced by a companion synthetic-data analysis and are
included via \texttt{\textbackslash figslot}, which falls back to a labelled
placeholder when the PDF is absent. The document therefore builds with or
without them, and no argument in the text depends on a figure's values.

\subsection{Limits of these computations}

Three caveats a reviewer should hold onto. The power calculations assume
bivariate normality and homoscedasticity; mastery pace is a rate and is likely
right-skewed, so the analysis should either transform it or use a model
appropriate to counts, and the sample sizes here should be treated as lower
bounds. The noncentrality convention is $\lambda=f^{2}n$, matching G*Power and
verified against its worked examples; other conventions give slightly different
answers and the difference is not negligible at small $n$. And every figure
depending on the screener's reliability inherits the status of that number,
which is a design target rather than a measurement.

A fourth caveat is more consequential than the other three and is stated in the
main text as well. The two parameters that most strongly drive the conclusions
--- the correlation between standing ability and learning rate, and the
comparator's reliability --- are unmeasured working values, not estimates. They
are carried through every table that depends on them as explicit sensitivity
ranges rather than points, and the computations are structured so that revising
either is a one-line change in \path{scripts/compute.py} followed by a rebuild.
A reader who disputes either value can therefore recompute the report's numbers
under their own assumption rather than arguing about the conclusion.

% ==========================================================================
\section*{References}
\addcontentsline{toc}{section}{References}
\small
\sloppy

\begin{itemize}[leftmargin=1.2em,itemsep=0.25em,label={}]

\item AERA, APA, \& NCME (2014). \textit{Standards for Educational and
Psychological Testing}. American Educational Research Association.

\item Cohen, J. (1988). \textit{Statistical Power Analysis for the Behavioral
Sciences} (2nd ed.). Lawrence Erlbaum.

\item Dixon, C., Oxley, E., Gellert, A. S., \& Nash, H. (2023). Muddying the
waters: A systematic review of dynamic assessment of reading.
\textit{Reading and Writing}, 36(3), 673--698.
\href{https://doi.org/10.1007/s11145-022-10312-3}{doi:10.1007/s11145-022-10312-3}

\item Dorans, N. J. (2004). Equating, concordance, and expectation.
\textit{Applied Psychological Measurement}, 28(4), 227--246.

\item Faul, F., Erdfelder, E., Buchner, A., \& Lang, A.-G. (2009). Statistical
power analyses using G*Power 3.1. \textit{Behavior Research Methods}, 41,
1149--1160.

\item Gelman, A., \& Stern, H. (2006). The difference between
``significant'' and ``not significant'' is not itself statistically
significant. \textit{The American Statistician}, 60(4), 328--331.
\href{https://doi.org/10.1198/000313006X152649}{doi:10.1198/000313006X152649}

\item Haberman, S. J. (2008). When can subscores have value?
\textit{Journal of Educational and Behavioral Statistics}, 33(2), 204--229.
\href{https://doi.org/10.3102/1076998607302636}{doi:10.3102/1076998607302636}

\item Haberman, S. J. (2008). Subscores and validity.
\textit{ETS Research Report Series}, 2008(2).
\href{https://doi.org/10.1002/j.2333-8504.2008.tb02150.x}{doi:10.1002/j.2333-8504.2008.tb02150.x}

\item Hanley, J. A., \& McNeil, B. J. (1982). The meaning and use of the
area under a receiver operating characteristic (ROC) curve.
\textit{Radiology}, 143(1), 29--36.
\href{https://doi.org/10.1148/radiology.143.1.7063747}{doi:10.1148/radiology.143.1.7063747}

\item Hanley, J. A., \& McNeil, B. J. (1983). A method of comparing the areas
under receiver operating characteristic curves derived from the same cases.
\textit{Radiology}, 148(3), 839--843.
\href{https://doi.org/10.1148/radiology.148.3.6878708}{doi:10.1148/radiology.148.3.6878708}

\item Holland, P. W., \& Dorans, N. J. (2006). Linking and equating. In R. L.
Brennan (Ed.), \textit{Educational Measurement} (4th ed., pp. 187--220).

\item Kane, M. T. (2013). Validating the interpretations and uses of test
scores. \textit{Journal of Educational Measurement}, 50(1), 1--73.

\item Killip, S., Mahfoud, Z., \& Pearce, K. (2004). What is an
intracluster correlation coefficient? Crucial concepts for primary care
researchers. \textit{Annals of Family Medicine}, 2(3), 204--208.
\href{https://doi.org/10.1370/afm.141}{doi:10.1370/afm.141}

\item Livingston, S. A., \& Lewis, C. (1995). Estimating the consistency and
accuracy of classifications based on test scores.
\textit{Journal of Educational Measurement}, 32(2), 179--197.

\item MacCallum, R. C., Zhang, S., Preacher, K. J., \& Rucker, D. D. (2002).
On the practice of dichotomization of quantitative variables.
\textit{Psychological Methods}, 7(1), 19--40.
\href{https://doi.org/10.1037/1082-989X.7.1.19}{doi:10.1037/1082-989X.7.1.19}

\item Messick, S. (1995). Validity of psychological assessment.
\textit{American Psychologist}, 50(9), 741--749.
\href{https://doi.org/10.1037/0003-066X.50.9.741}{doi:10.1037/0003-066X.50.9.741}

\item Lohman, D. F. (2005). An aptitude perspective on talent identification.
\textit{Gifted Child Quarterly}, 49(2), 111--138.
\href{https://doi.org/10.1177/001698620504900203}{doi:10.1177/001698620504900203}

\item Meng, X.-L., Rosenthal, R., \& Rubin, D. B. (1992). Comparing correlated
correlation coefficients. \textit{Psychological Bulletin}, 111(1), 172--175.

\item NWEA (2025). \textit{MAP Growth Technical Report}. NWEA.

\item Ozen, Z., et al. (2024). A meta-analysis of the criterion-related validity
of the Cognitive Abilities Test. \textit{Gifted Child Quarterly}.
\href{https://doi.org/10.1177/00169862241285593}{doi:10.1177/00169862241285593}

\item Peters, S. J., Rambo-Hernandez, K., Makel, M. C., Matthews, M. S., \&
Plucker, J. A. (2019). Effect of local norms on racial and ethnic representation
in gifted education. \textit{AERA Open}, 5(2).

\item Sackett, P. R., Zhang, C., Berry, C. M., \& Lievens, F. (2022).
Revisiting meta-analytic estimates of validity in personnel selection:
Addressing systematic overcorrection for restriction of range.
\textit{Journal of Applied Psychology}, 107(11), 2040--2068.
\href{https://doi.org/10.1037/apl0000994}{doi:10.1037/apl0000994}

\item Sinharay, S. (2010). When do subscores have added value? Criteria based on
classical test theory. \textit{Journal of Educational Measurement}, 47(2),
150--174.

\item Taylor, H. C., \& Russell, J. T. (1939). The relationship of validity
coefficients to the practical effectiveness of tests in selection.
\textit{Journal of Applied Psychology}, 23(5), 565--578.

\item Thorndike, R. L. (1949). \textit{Personnel Selection: Test and Measurement
Techniques}. Wiley.

\item VanderWeele, T. J., \& Ding, P. (2017). Sensitivity analysis in
observational research: Introducing the E-value.
\textit{Annals of Internal Medicine}, 167(4), 268--274.

\item Westfall, J., \& Yarkoni, T. (2016). Statistically controlling for
confounding constructs is harder than you think.
\textit{PLoS ONE}, 11(3), e0152719.
\href{https://doi.org/10.1371/journal.pone.0152719}{doi:10.1371/journal.pone.0152719}

\end{itemize}

\normalsize

\subsection*{Repository sources}

Claims about the instrument's current state are grounded in the project's own
records rather than external literature:
\path{docs/product/STAGE2_BANK_RECOVERY_MEASUREMENT.md} (recovery ladder,
contamination floor, indeterminate rates);
\path{docs/product/STAGE2_QUESTION_DESIGN.md} (block design and gates);
\path{docs/research/ASSUMPTIONS_AND_EVIDENCE.md} entries E-095, E-100,
E-101, E-200 (recovery, validation-data commitment, cohort sizes, estimator
defect);
\path{docs/governance/DECISION_LOG.md} entries D-030 and D-200 (learning-rate
reportability);
\path{docs/research/IN_HOUSE_COGNITIVE_TEST_RESEARCH.md} and
\path{docs/research/EVIDENCE_DOSSIER.md} (CogAT and MAP psychometric
properties);
\path{docs/interviews/CRYSTAL_MARTEL_INTERVIEW_NOTES.md} (cohort sizes,
seat rationing, validation-data availability).
